\subsection{Published-Solver Comparison and Joint-Shift Stress}
\label{sec:v06_performance}

We evaluate all 23 frozen policy identities on 302 contexts at nominal wall
 targets $T\in\{0.1,1,5\}$ seconds. Newly drawn scheduling problems provide
216 endpoints from 108 resource-intervention pairs. Public inputs comprise
25 WDP, three Segmentation, and ten Grids sources from the weighted-clique
benchmark~\cite{mccreesh2017evaluation}, retaining exact weights and complement
orientation. Public source clusters follow fixed withholding within previously
exposed corpora. Two 24-context C3 tracks remain separate exploratory analyses.

Table~\ref{tab:v06-published} displays four predeclared families at $T=5$
against CHILS and ILS~\cite{grossmann2025chils}, M2WIS~\cite{grossmann2024mmwis},
Struction~\cite{gellner2021struction}, and WeightedBR~\cite{lamm2019weighted}.
The deployment comparison keeps the same strong CHILS incumbent and repair
limits for joint witness, Degree, and EoH-DSL policies. Four frozen EoH-DSL
outputs retain their common seed AST~\cite{liu2024eoh}. Every warm policy is
charged the full shared initializer and its own repair cost; nominal targets
are not equal end-to-end deadlines across C++ and compiled Python.

Fresh contact pairs change ground gap $0.5\!\to\!6$ with unchanged
opportunities. They also increase size to 512/1,024/2,048 contacts and vary
weights, durations and density. This combined shift tests deployment stress;
it cannot isolate the effect of unseen gaps. Arbitrary public graphs test
solver compatibility, not satellite constraint adaptation.

Figure~\ref{fig:v06_performance_budget} plots measured CPU against size at
fixed $T=5$ for two predeclared contact families. Points retain paired-source
costs; lines join their means. Complete cold/warm reward, coverage and all
13-family/three-target matrices remain in the supporting report. Cohort
reward means require every member to return a successful feasible schedule;
partial means cannot establish superiority over fully covered comparators.
No test outcome changes a program, budget or displayed family.

The common optimization backbone retains its feasible incumbent while the
synthesized policies provide refined decision representations. The deployment
results illustrate this separation between information adequacy and achieved
schedule quality; they do not establish an incremental LLM quality gain.
